% Seccion 6 de la rubrica. Cada tecnica ya reporta e interpreta sus propios
% resultados en la seccion 5; aqui se consolidan en una sola tabla para que
% sean comparables, sin repetir la interpretacion.

\section{Resultados}
\label{sec:resultados}

La \cref{tab:resultados} consolida el desempeño de los tres clasificadores
sobre la validación de 2025 junto con las dos líneas base del modelo
autorregresivo, y la \cref{tab:resultados-ejes} resume lo que cada técnica
de interdependencia entregó. Los tres clasificadores superan el criterio de
logro de 0.75 fijado en \cref{sec:obj:objetivo}.

\begin{table}[t]
  \centering
  \caption{Desempeño en la validación de 2025. Los modelos día-estación
    (logística y discriminante) y el modelo de área metropolitana
    (autorregresivo y sus líneas base) responden preguntas distintas y sus
    AUC no son directamente comparables entre bloques. Sensibilidad y
    especificidad al umbral de Youden fijado en el ajuste; la
    generalización es el AUC fuera de muestra por validación cruzada
    dejando un año fuera.}
  \label{tab:resultados}
  \small
  \begin{tabular}{lrrrr}
    \toprule
    Modelo & AUC & Sens. & Esp. & General. \\
    \midrule
    \multicolumn{5}{l}{\emph{Día-estación} ($n=2{,}927$)} \\
    Regresión logística    & 0.850 & 0.698 & 0.853 & 0.80 \\
    Discriminante lineal   & 0.846 & 0.758 & 0.788 & 0.80 \\
    \midrule
    \multicolumn{5}{l}{\emph{Área metropolitana} ($n=365$)} \\
    Autorregresivo AR(1)   & 0.898 & 0.759 & 0.841 & 0.878 \\
    Climatología (base)    & 0.668 & ---   & ---   & --- \\
    Persistencia (base)    & 0.717 & ---   & ---   & --- \\
    \bottomrule
  \end{tabular}
\end{table}

\begin{table}[t]
  \centering
  \caption{Ejes retenidos por las técnicas de interdependencia sobre el
    ajuste 2021--2024. Varianza: total en el ACP, común en el AF.}
  \label{tab:resultados-ejes}
  \small
  \begin{tabular}{llr}
    \toprule
    Técnica & Eje & Var. (\%) \\
    \midrule
    \multirow{3}{*}{ACP (15 variables)}
      & Carga primaria de combustión & 26.1 \\
      & Forzamiento fotoquímico      & 19.8 \\
      & Estancamiento atmosférico    & 10.0 \\
    \midrule
    \multirow{3}{*}{AF (13 variables)}
      & Combustión primaria          & 15.7 \\
      & Ventilación                  & 14.5 \\
      & Forzamiento fotoquímico      & 10.1 \\
    \bottomrule
  \end{tabular}
\end{table}

Cuatro resultados atraviesan las técnicas. Primero, los mismos tres ejes
aparecen por dos caminos distintos, y el AF corrige al ACP en un punto: el
eje de estancamiento era un eje de viento, porque la presión barométrica
tiene comunalidad de 0.13 y casi toda su varianza es específica. Segundo, la
jerarquía de efectos es estable en los tres clasificadores: el forzamiento
fotoquímico pesa más que la combustión primaria, la ventilación protege, y
en los modelos día-estación la temporada y la zona pesan más que cualquier
factor. Tercero, la dependencia temporal es real y aprovechable: el rezago
de la respuesta sube el AUC de área metropolitana de 0.882 a 0.898 y
multiplica los momios por 3.83. Cuarto, todos los modelos comparten un
efecto de año que ningún predictor explica: 2024 y 2025 son los años de
ozono más alto de la serie y también los mejor clasificados, y el desempeño
esperado en un año cualquiera queda entre 0.04 y 0.05 por debajo del de la
validación.
